% Generated by scripts/assemble.py; corrections belong in source/corrections.json.
\examhead{28}{中国石油大学(华东) 2024 年数学分析试卷}{32}
% p033-b01; source PDF page 33
\par\Needspace{6\baselineskip}
\blockanchor{p033-b01}
\question{1}(10 分) 求极限\par
\nopagebreak[4]
% p033-b02; source PDF page 33
\blockanchor{p033-b02}
\displayformula{\lim _ {n \rightarrow \infty} \left(\frac {\pi}{\sqrt {n ^ {2} - 1}} + \frac {\pi}{\sqrt {n ^ {2} - 2}} + \dots + \frac {\pi}{\sqrt {n ^ {2} - 2 n}}\right).}
% p033-b03; source PDF page 33
\blockanchor{p033-b03}
\question{2}(10 分) 求极限 $\lim_{x \to 0} \left( \frac{\sin x}{x} \right)^{\frac{1}{x^2}}$.\par
% p033-b04; source PDF page 33
\blockanchor{p033-b04}
\question{3}(10 分) 对任意的自然数 $n$, 有 $a_{n} \leq b_{n} \leq c_{n}$, 级数 $\sum_{n=1}^{\infty} a_{n}$ 与 $\sum_{n=1}^{\infty} c_{n}$ 均收敛, 证明: 级数 $\sum_{n=1}^{\infty} b_{n}$ 也收敛.\par
% p033-b05; source PDF page 33
\par\Needspace{6\baselineskip}
\blockanchor{p033-b05}
\question{4}(10 分) 设 $f(x)$ 在 $[0,1]$ 上连续, 且 $f(\sin x) + f(\cos x) \neq 0$, 求积分\par
\nopagebreak[4]
% p033-b06; source PDF page 33
\blockanchor{p033-b06}
\displayformula{\int_ {0} ^ {\frac {\pi}{2}} \frac {f (\sin x)}{f (\sin x) + f (\cos x)} d x.}
% p033-b07; source PDF page 33
\blockanchor{p033-b07}
\question{5}(10 分) 求 $x^{2} + (y - b)^{2} = a^{2}$ ($b > a > 0$) 绕 $x$ 轴旋转所得的旋转体的体积.\par
% p033-b08; source PDF page 33
\par\Needspace{6\baselineskip}
\blockanchor{p033-b08}
\question{6}(10 分) 设数列 $\{a_{n}\}$ 满足: 存在正数 $M > 0$, 对一切 $n$ 有\par
\nopagebreak[4]
% p033-b09; source PDF page 33
\blockanchor{p033-b09}
\displayformula{A _ {n} = \left| a _ {2} - a _ {1} \right| + \left| a _ {3} - a _ {2} \right| + \dots + \left| a _ {n} - a _ {n - 1} \right| \leq M.}
% p033-b10; source PDF page 33
\blockanchor{p033-b10}
证明: 数列  $\{a_{n}\}$  与  $\{A_{n}\}$  都收敛.\par
% p033-b11; source PDF page 33
\par\Needspace{6\baselineskip}
\blockanchor{p033-b11}
\question{7}(10 分) 记\par
\nopagebreak[4]
% p033-b12; source PDF page 33
\blockanchor{p033-b12}
\displayformula{F (t) = \iiint_ {V} x \ln \left(1 + x ^ {2} + y ^ {2} + z ^ {2}\right) \mathrm{d} x \mathrm{d} y \mathrm{d} z.}
% p033-b13; source PDF page 33
\blockanchor{p033-b13}
其中 $V: x^{2} + y^{2} + z^{2} \leq t, \sqrt{y^{2} + z^{2}} \leq x,$ 求 $\frac{\mathrm{d}F(t)}{\mathrm{d}t}$.\par
% p033-b14; source PDF page 33
\blockanchor{p033-b14}
\question{8}(10 分) 设 $\Omega$ 是由 $z = \sqrt{9 - x^2 - y^2}$ 与 $z = \sqrt{x^2 + y^2}$ 围成的区域, $f(x,y,z)$ 为 $\Omega$ 上的连续函数, 将 $\iiint_{\Omega} f(x,y,z) \mathrm{d}x \mathrm{d}y \mathrm{d}z$ 分别用直角坐标系, 柱面坐标, 球面坐标化为三重积分(无需计算).\par
% p033-b15; source PDF page 33
\blockanchor{p033-b15}
\question{9}(10 分) 设 $f(x)$ 在 $x = 0$ 的某邻域内存在连续的二阶导数, 且 $\lim_{x \to 0} \frac{f(x)}{x} = 0$. 证明: 级数 $\sum_{n=1}^{\infty} f\left(\frac{1}{n}\right)$ 绝对收敛.\par
% p033-b16; source PDF page 33
\par\Needspace{6\baselineskip}
\blockanchor{p033-b16}
\question{10}(10 分) 计算曲面积分\par
\nopagebreak[4]
% p033-b17; source PDF page 33
\blockanchor{p033-b17}
\displayformula{\iint_ {S} (x - x y) \mathrm{d} y \mathrm{d} z + (y - y z) \mathrm{d} z \mathrm{d} x + (z - x y) \mathrm{d} x \mathrm{d} y.}
% p033-b18; source PDF page 33
\blockanchor{p033-b18}
其中 S 是以原点为球心, 半径为 a 的上半球面, 方向取上侧.\par
% p033-b19; source PDF page 33
\blockanchor{p033-b19}
\question{11}(10 分) 求 $y = \ln (1 - 3x + x^2)$ 的麦克劳林级数展开式.\par
% p033-b20; source PDF page 33
\blockanchor{p033-b20}
\question{12}(10 分) 设 $f(x)$ 在 $(- \infty, + \infty)$ 上连续, 且 $\lim_{x \to \infty} \frac{f(x)}{x} = 0$, 证明: 存在 $\xi \in (-\infty, + \infty)$, 使得 $f(\xi) = \xi$.\par
% p033-b21; source PDF page 33
\blockanchor{p033-b21}
\question{13}(10 分) 设 $f(x, y) = x + y + xy$, 曲线 $C: x^2 + y^2 = 2$.\par
% p033-b22; source PDF page 33
\blockanchor{p033-b22}
(1) 求 $f(x)$ 在 $P(1,1)$ 沿方向 $l = (-1, -1)$ 的方向导数.\par
% p033-b23; source PDF page 33
\blockanchor{p033-b23}
(2) 求  $f(x)$  在曲线 C 上的最大方向导数.\par
% p033-b24; source PDF page 33
\blockanchor{p033-b24}
\question{14}(10 分) 设 $f(x)$ 在 $[1,2]$ 上可导, 且 $f'(x) \neq 1, \forall x \in [1,2]$. 若对任意的 $x \in [1,2]$, 有 $f(x) \in (1,2)$, 证明: 存在唯一的 $\xi \in [1,2]$, 使得 $f(\xi) = \xi$.\par
% p033-b25; source PDF page 33
\blockanchor{p033-b25}
\question{15}(10 分) 设 $f(x)$ 在 $[a, b]$ 上可积, 证明: $f^2(x)$ 在 $[a, b]$ 上也可积.\par
